% ═══════════════════════════════════════════════════════════════════
% Chapter 10 — Mimic: Per-Vendor Backend Framework
% ═══════════════════════════════════════════════════════════════════

\section{Mimic: Per-Vendor Backend Framework}
\label{sec:mimic}

Mimic is Crucible's per-vendor backend framework.
It consumes \IR{2} (Chapter~\ref{sec:ir002}) and produces vendor binary code via per-vendor implementations of \IRstar{3} (Chapter~\ref{sec:ir003}).
Mimic is structurally a shared core plus one self-contained backend per vendor: \code{mimic/nv/} for NVIDIA, \code{mimic/am/} for AMD, \code{mimic/tpu/} for Google, \code{mimic/trn/} for AWS Trainium, \code{mimic/cer/} for Cerebras, and \code{mimic/cpu/} for x86, AArch64, and RISC-V.

Each backend owns its vendor's \IRstar{3} representation, its instruction encoder and decoder, its three-tier simulator with calibrated cost model, its MAP-Elites kernel search, its runtime library, its collective library, and its calibration harness.
The shared core provides the simulator framework, MAP-Elites driver, archive machinery, mutation dispatch, the cross-vendor numerics CI harness, calibration skeleton, and effect-token plumbing.

Mimic has no vendor-SDK dependency: zero cuBLAS, zero cuDNN, zero NCCL, zero rocBLAS, zero MIOpen, zero RCCL, zero libtpu, zero libnrt, zero libnccom, zero hcoll, zero libsharp.
Each per-vendor runtime library wraps the kernel-driver ioctl protocol directly; each per-vendor collective library composes over the native fabric primitives via Crucible Native Transport Protocol (Chapter~\ref{sec:cntp}).
The kernel driver (NVIDIA's \code{nvidia.ko}, AMD's \code{amdkfd}, Google's \code{accel\_*}, AWS's \code{neuron.ko}) is the only vendor-proprietary surface Crucible inherits.

This chapter specifies Mimic's architecture: the per-vendor template, the public-API surface that Forge calls, the simulator-and-search internals, and the cross-vendor calibration discipline.
The per-vendor specifications --- the SASS encoder for Hopper and Blackwell, the AMDGPU emitter for CDNA3 and RDNA3, the TPU executable format, the NEFF format, the Cerebras Software Language emitter --- are documented in the misc/MIMIC.md design specification.

\subsection{The per-vendor template}
\label{sec:mimic-template}

Every per-vendor backend implements ten deliverables.
The template is uniform across vendors; adding a new vendor is the cost of implementing the template once, with no changes to Forge or to the shared core.

\begin{enumerate}
  \item \textbf{\IRstar{3} type system, builder, verifier.}
        The vendor-specific machine intermediate representation, encoded with the same axiomatic discipline as \IR{1} and \IR{2}: arena-owned, NSDMI on every field, strong-typed identifiers.
  \item \textbf{\IR{2}-to-\IRstar{3} lowering.}
        Conventional Static Single Assignment construction, address-space resolution, register allocation, instruction scheduling, recipe realization (Section~\ref{sec:mimic-recipe-realization}).
  \item \textbf{Encoder and binary-format writer.}
        \IRstar{3} to ISA bytes; ISA bytes to vendor binary format (cubin on NVIDIA, HSACO on AMD, executable on TPU, NEFF on Trainium, CSL on Cerebras, native ELF on CPU).
  \item \textbf{Decoder.}
        Inverse of the encoder, with round-trip correctness (Section~\ref{sec:mimic-roundtrip}).
        Used for calibration validation, debugging, and loading cached binaries from Cipher (Chapter~\ref{sec:cipher}).
  \item \textbf{Three-tier simulator.}
        Fast tier (interval-based, 1 to 5 milliseconds per simulation, ranking-accurate), medium tier (event-driven, 10 to 30 milliseconds, 95-percent ranking accuracy and 5-to-8-percent absolute), accurate tier (cycle-stepping, 100 to 500 milliseconds, 99-percent ranking and 2-to-5-percent absolute).
        Each tier shares the simulator framework from the shared core; per-vendor specialization adds the memory subsystem, the warp or wave scheduler, and the async-pipeline models.
  \item \textbf{MAP-Elites mutation operators and insight catalog.}
        The shared core ships approximately twenty vendor-agnostic insight kinds (high register pressure, fast-local-storage saturation, queue contention, scheduler imbalance); each backend adds approximately ten vendor-specific kinds (\code{WGMMA\_UNDERUTILIZED} on NVIDIA, \code{MFMA\_WAVE\_STALL} on AMD, \code{MXU\_INPUT\_STARVATION} on TPU).
        Mutation operators are similarly factored into vendor-agnostic and vendor-specific.
  \item \textbf{Peephole rule set.}
        Approximately one hundred fifty rules per vendor, table-driven, applied in a single pass after register allocation and before encoding.
  \item \textbf{Calibration harness.}
        Approximately one hundred fifty microbenchmarks per vendor, run once per chip target, producing the calibrated \struct{TargetCaps} and \struct{OpcodeLatencyTable} (Section~\ref{sec:mimic-calibration}).
  \item \textbf{Runtime library.}
        Replacement for libcuda, libnrt, libtpu, libhsa, with the same surface (device discovery and open, memory pool management, binary loading, kernel submission, completion polling, graph capture).
        Wraps the kernel-driver ioctl protocol directly.
  \item \textbf{Collective library.}
        Replacement for NCCL, RCCL, libnccom, hcoll, with the same surface (all-reduce, all-gather, reduce-scatter, all-to-all, broadcast, point-to-point) and additional CNTP integration for cross-vendor and cross-node fabrics.
\end{enumerate}

The shared core (\code{mimic/core/}) provides: the simulator framework skeleton (event queue, scheduler skeleton, memory-model skeleton), the MAP-Elites driver and Archive structure, the mutation dispatcher, the base insight catalog, the calibration-harness framework, the cross-vendor numerics CI harness (Chapter~\ref{sec:cross-vendor-ci}), and the standard arena and effect-token infrastructure.
Each backend specializes the framework with its own memory subsystem, scheduler, and async pipelines.

\subsection{The public interface}
\label{sec:mimic-public-interface}

Mimic exposes five entry points to Forge.
The signatures are uniform across vendors; the dispatch by \code{TargetCaps::vendor\_id} is internal to Mimic's top-level shim.
The five entries cover the full Forge--Mimic interaction surface; nothing else crosses the boundary.

\begin{enumerate}
  \item \code{mimic::fast\_cost(KernelNode, TargetCaps) -> Cycles}.
        Approximate cycle count from the fast tier of the per-vendor simulator.
        Approximately ten microseconds per call.
        Used by Forge's Phase~B.5 (per-node cost estimation), Phase~D.2 (fusion-cost evaluation), and Phase~F.2 (tile pruning).
        The estimate is ranking-accurate but not absolute-accurate; ten-to-twenty-percent absolute error is admissible as long as it preserves the ordering between candidates.

  \item \code{mimic::propose\_tiles(KernelNode, TargetCaps) -> Span<AbstractTile>}.
        Generates abstract tile candidates that are hardware-realizable on the current target.
        Used by Forge's Phase~F.1 (tile proposal).
        Sixteen to sixty-four candidates per call; per-vendor pruning of obviously-infeasible candidates is part of the implementation.

  \item \code{mimic::compile\_kernel(KernelNode, TargetCaps, CompileConfig) -> CompiledKernel}.
        The principal entry: produces an opaque \struct{CompiledKernel} with vendor binary bytes, predicted cycles, and structured insights.
        Used by Forge's Phase~H.
        Tens to hundreds of milliseconds per kernel; runs MAP-Elites internally; parallelizable across kernels via Forge's thread pool.
        The \struct{CompileConfig} parameter selects the search mode (MAP-Elites only, MAP-Elites with hardware validation, or beam search) and the wall-clock and iteration budgets.

  \item \code{mimic::predict(CompiledKernel, TargetCaps, SimTier) -> SimResult}.
        Runs the simulator on an already-compiled kernel without re-compiling.
        Used by Forge's Phase~L.2 (drift detection): the predicted cycles are compared against the measured cycles for a sampled invocation.
        The simulator tier is configurable (fast for cheap drift checks, medium for primary use, accurate for calibration validation).

  \item \code{mimic::probe\_counters(CompiledKernel, TargetCaps) -> Measurements}.
        Reads the per-vendor hardware counters via the native profiling interface (CUPTI on NVIDIA, rocprof on AMD, the PJRT profiler on TPU, neuron-profile on Trainium, perf\_event\_open on CPU).
        Returns a uniform \struct{Measurements} record carrying twelve signals.
        Used by Forge's Phase~L (continuous validation).
\end{enumerate}

The dispatch shim is mechanical: the top-level \code{mimic::compile\_kernel} (and analogously each entry) reads \code{caps.vendor\_id} and delegates to the per-vendor namespace.
The cost is a jump table; there is no runtime polymorphism, no virtual dispatch, no allocation.

\subsection{The three subsystems}
\label{sec:mimic-three-subsystems}

Each per-vendor backend factors into three loosely coupled subsystems sharing target-model data: an encoder--decoder pair, a simulator (three tiers), and a MAP-Elites synthesizer.
The shared data is the calibrated \struct{TargetCaps} and the \struct{OpcodeLatencyTable} (Section~\ref{sec:mimic-target-model}); each subsystem is independently testable.

The encoder--decoder pair is publicly callable: Forge's Phase~J (\textsc{emit}, Section~\ref{sec:forge-phase-j}) calls the encoder; Augur's recovery path (Chapter~\ref{sec:augur}) calls the decoder.
Round-trip correctness is asserted continuously: \code{decode(encode(program))} equals \code{program} for every valid \IRstar{3} program.

The simulator is publicly callable through \code{mimic::predict}.
It also runs internally during MAP-Elites search.
The simulator tiers form an escalation chain (Section~\ref{sec:mimic-simulator}).

The MAP-Elites synthesizer is the implementation of \code{mimic::compile\_kernel}.
It uses the encoder to produce candidates, the simulator to score them, the archive to maintain diversity, and the decoder to verify round-trip correctness on the chosen winner.

\subsection{The target model}
\label{sec:mimic-target-model}

Each per-vendor backend exposes a \struct{TargetCaps} structure that encodes the calibrated capabilities of a chip.
The shared core defines the \emph{abstract} fields that Forge consumes; each backend extends with vendor-specific fields opaque to Forge.

\begin{definition}[Abstract TargetCaps]
\label{def:mimic-targetcaps}
The abstract \struct{TargetCaps} structure carries: identity (vendor enumeration, chip family, chip variant, suffix flag); compute geometry (execution-unit count, threads per execution unit, maximum registers per thread, maximum concurrent threads per unit); fast-local-storage capacity; tensor-fragment-storage capacity; cacheable-storage capacity; main-memory capacity; bandwidth estimates (main memory, L2 cache, fast local, intra-node interconnect, inter-node interconnect); launch-overhead estimate; and an opaque pointer to the per-vendor extension.
The structure is one cache line; vendor-specific fields live in a separately allocated \struct{TargetCaps<Vendor>} pointed at by the opaque field.
\end{definition}

The split between abstract and per-vendor fields is the structural mechanism that keeps Forge vendor-neutral.
Forge reads \code{TargetCaps::execution\_units} and \code{TargetCaps::fast\_local\_bytes}; it never reads the per-vendor extension.
The opaque pointer is dereferenced only by Mimic's per-vendor functions.

The \struct{OpcodeLatencyTable} is per-SM (or per-CU, per-MXU) data: the latency, occupancy, throughput-inverse, and pipeline-assignment of every operator the per-vendor backend emits.
The table is approximately one hundred kilobytes per chip target, loaded lazily from a JSON data file on first use.
The data files (\code{crucible/mimic/<vendor>/data/<chip>.json}) are checked into version control and regenerated by the calibration harness.

\subsection{The three-tier simulator}
\label{sec:mimic-simulator}

The simulator is the principal source of cost information in Mimic.
A single simulator cannot serve both MAP-Elites throughput and calibration-validation accuracy; the answer is three tiers sharing model data.

\textbf{Fast tier.}
Interval-based critical-path analysis with analytical memory-bandwidth model and queue-depth estimation.
Wall time is approximately one to five milliseconds per simulation on CPU, ten times faster on GPU when batched.
Throughput is ten thousand candidates per second per CPU core or one hundred thousand per second on GPU.
Absolute accuracy is ten to twenty percent; ranking accuracy is approximately ninety percent.
The fast tier skips cycle-level warp scheduling, register-bank conflicts, operand-collector detail, the DRAM row-buffer model, and cache state.
Used for: MAP-Elites coarse filtering (most candidates are rejected at the fast tier without further evaluation).

\textbf{Medium tier.}
Event-driven warp-scheduler simulator with scoreboard tracking, explicit memory-subsystem queues (L1, L2, main memory), functional-unit token pools, operand collector, per-warp state machines.
Wall time is approximately ten to thirty milliseconds.
Throughput is one hundred to three hundred candidates per second per CPU core.
Absolute accuracy is five to eight percent; ranking accuracy exceeds ninety-five percent.
Used for: MAP-Elites primary fitness evaluation, Forge's Phase~B.5 cost estimates, and Augur's prediction.

\textbf{Accurate tier.}
Cycle-by-cycle simulation of every pipe, register read port, scoreboard slot, DRAM command, bank-conflict resolution, and commit-group-aware async tracking.
Wall time is approximately one hundred to five hundred milliseconds.
Throughput is two to ten candidates per second.
Absolute accuracy is two to five percent; ranking accuracy exceeds ninety-nine percent.
Used for: calibration validation (the harness compares accurate-tier predictions against measured hardware on a regression set of fifty kernels), MAP-Elites tie-breaks on top candidates per generation, and debugging simulator discrepancies.

The escalation policy in MAP-Elites is structural: every candidate is fast-tier evaluated; the bottom seventy percent are rejected; the top thirty percent receive medium-tier scoring; the top ten by medium fitness receive accurate-tier evaluation; the archive update uses the most accurate tier that evaluated the candidate.

\subsection{Memory subsystem modeling}
\label{sec:mimic-memory-model}

The memory subsystem is the hardest part of GPU simulation and the principal source of accuracy loss in academic simulators.
Mimic resolves the difficulty by leveraging the open NVIDIA driver source (through Hopper) and AMD's published register specifications: the memory-subsystem mechanisms in Mimic are derived from documented driver behavior, not reverse-engineered from microbenchmark residuals.

The directly documented mechanisms include the L2 slice address-hash function (the GPC-to-FBP crossbar routing algorithm), the DRAM command scheduler (FR-FCFS with per-channel age prioritization), the row-buffer policy (open-page by default, with closing trigger on hit-rate drop), per-channel refresh scheduling and stall propagation, the L1 eviction policy (LRU with replacement bypass for streaming operations), the coalescer logic (128-byte transaction merging), the read-write turnaround budget, and the error-correction overhead percentage.

The medium-tier model represents these mechanisms as a queue-aware latency function: an L1 hit returns at the calibrated L1 latency; an L1 miss routes to an L2 slice via the documented hash; the L2 hit latency is the base latency plus the queue-depth multiplier; an L2 miss routes to a main-memory channel; the main-memory latency is the row-hit or row-miss latency plus the read-write-turnaround penalty plus the queue-depth contribution.
Refresh stalls are scheduled per channel.

The accurate tier represents the mechanisms in full: per-bank row-buffer state explicit, DRAM commands scheduled individually, per-slice and per-channel queues simulated, refresh interleaved with reads and writes, write-combining buffers tracked, L2 hardware compression ratio applied.
The accurate tier is approximately twenty times slower than the medium tier in exchange for two-to-three additional percentage points of absolute accuracy.

The remaining residual error (two to five percent at the accurate tier) is attributable to effects below the silicon-software boundary: temperature-dependent DRAM timing, cross-kernel cache retention, exact compression ratio under non-microbenchmark workloads, and crossbar-arbitration details under extreme contention.
These are not fixable from the driver source; they are accepted as the accuracy floor.

\subsection{The warp scheduler model}
\label{sec:mimic-warp-scheduler}

NVIDIA's warp scheduler is Greedy-Then-Oldest with age-based boosting.
The scheduler-level details are not publicly documented but are accessible from the driver source through Hopper.
Mimic represents the policy as three interleaved rules: the greedy rule (stick with the last-picked warp while it is ready and within the burst budget), the oldest rule (fall back to the oldest-ready warp when greedy fails), and the boost rule (elevate any warp that has been stuck beyond the boost threshold).

The three constants --- the greedy probability, the boost threshold, and the burst limit --- are calibrated via dedicated microbenchmarks: a two-warp benchmark with one fast and one slow warp fits the greedy probability; an N-warp benchmark with one stalled warp fits the boost threshold; a single-warp benchmark with a long dependency chain fits the burst limit.
The same calibration discipline applies, with vendor-specific instantiations, to AMD's wave scheduler and TPU's scalar-processor dispatcher.

\subsection{Async-pipeline modeling}
\label{sec:mimic-async}

NVIDIA Hopper introduces three async subsystems: WGMMA (warpgroup MMA), TMA (Tensor Memory Accelerator), and (on Blackwell) tcgen05 with TMEM.
Each is modeled in Mimic's medium and accurate tiers as a separate pool of outstanding operations with explicit issue, deadline, and completion events.

The principal correctness concern is the back-to-back pipeline overlap.
A WGMMA's MMA pipe is busy for the \code{occupancy} cycles after issue, but the result does not land until \code{latency} cycles after issue.
Other operations issue during the gap, delivering pipeline overlap.
Modeling correctness requires releasing the MMA pipe token at \code{issue\_time + occupancy}, not at \code{issue\_time + latency}.
Modeling this incorrectly underestimates throughput by a factor of two to eight.

The same discipline applies, vendor-specifically, to AMD's MFMA (no asynchronous interface; \code{s\_waitcnt}-based synchronization), TPU's MXU (scheduled by the scalar processor), and Trainium's TensorEngine (per-engine state).

\subsection{Insight extraction}
\label{sec:mimic-insights}

Insights are Mimic's second-most-important output after cycles.
An insight is a structured diagnostic produced during simulation: a kind enumeration, a severity, a confidence, an instruction-level localization, a metric value, and an optional mutation hint.
Approximately forty insight kinds are defined: high register pressure, fast-local-storage bank conflicts, main-memory bandwidth saturation, mid-pipeline contention, async-operation underutilization, scheduler imbalance, and so on.
Approximately twenty kinds are vendor-agnostic and live in the shared core; the remainder are per-vendor.

Each insight kind carries a canonical mutation hint mapping to one to three mutation operators.
For instance, \code{REGISTER\_PRESSURE\_HIGH} hints at \code{DECREASE\_TILE\_SIZE}, \code{ENABLE\_UR\_PROMOTION}, or \code{REDUCE\_PIPELINE\_STAGES}; \code{WGMMA\_UNDERUTILIZED} hints at \code{INCREASE\_PIPELINE\_STAGES}, \code{INCREASE\_TILE\_K}, or \code{ENABLE\_WARP\_SPECIALIZATION}.
The mapping is the structural mechanism that drives MAP-Elites's mutation selection: an insight-driven mutation is approximately three to eight times more effective than a random mutation, because it applies expert knowledge encoded in the catalog.

Insights are extracted during simulation, not as a post-hoc pass.
Each event handler in the simulator (instruction completion, scoreboard release, memory response) carries a hook that examines local state and emits an insight when a threshold is met.
The cost is approximately ten to fifty nanoseconds per hook.

\subsection{MAP-Elites kernel search}
\label{sec:mimic-map-elites}

MAP-Elites is Mimic's solution to the kernel-synthesis problem: given an \IR{2} \struct{KernelNode}, find an \IRstar{3} program that minimizes cycles subject to target constraints.
Unlike greedy search (XLA's benchmark-best-wins) or random autotuning (Triton's config sweep), MAP-Elites maintains a behavior archive: a discretized grid where each cell stores the best solution at that behavior coordinate.

\begin{definition}[Behavior archive]
\label{def:mimic-archive}
The MAP-Elites archive is a six-dimensional grid with eight buckets per dimension: occupancy, register usage, fast-local-storage usage, pipeline depth, MMA shape family, and warpgroup split.
Total cells: 8\textsuperscript{6} = approximately 260\,000.
A typical kernel family populates 500 to 5\,000 cells.
Each cell holds the best (highest-fitness) candidate seen at that behavior coordinate, plus the candidate's content hash, fitness, mutation count, and structured insights.
\end{definition}

The fitness function is the efficiency ratio: theoretical-minimum cycles divided by simulated cycles, normalized to the unit interval.
Perfect roofline utilization is 1.0; 0.5 is fifty-percent efficient; 0.9 is excellent; 0.95 is near-optimal; values above 0.98 are typically measurement error.

The mutation catalog has approximately fifty operators in three classes.
\textbf{Macro mutations} (about ten percent of applications) change the structural shape: tile-shape swap, warpgroup-split reassignment, grid-shape adjustment, cluster-size change, layout swap, split-K toggle, persistent-CTA conversion.
\textbf{Meso mutations} (about forty percent) change the pipeline structure without altering tiles: pipeline-stages adjustment, swizzle-pattern change, MMA-shape variant, epilogue-scheme change, prefetch-stage addition or removal, register-allocation strategy change.
\textbf{Micro mutations} (about fifty percent) make small tweaks: register-pair swap, instruction reordering, reuse-hint addition, warp-order adjustment, scoreboard-allocation change, stall-count adjustment.

Mutation selection is insight-driven: with seventy-percent probability the mutator samples an insight from the parent's simulation result and applies one of the canonical hints; with thirty-percent probability the mutator samples uniformly across the macro/meso/micro distribution.

The main loop runs concurrent batches of mutations against a thread pool sized by hardware concurrency, evaluates each batch through the simulator-tier escalation, and submits the candidates to the archive.
Termination conditions are: efficiency ratio above 0.92 (near-optimal; stop), wall-clock budget exceeded, iteration budget exceeded, or fifty generations of stuck-at-best (no improvement).

\textbf{Warm-start.}
The archive is persisted to Cipher between compilations.
A subsequent compilation of a structurally similar kernel (the same \struct{KernelKind}, comparable shape, comparable dtype) loads the previous archive and seeds MAP-Elites from the highest-fitness cells.
The archive lookup is two-key: an exact match on \struct{KernelContentHash} returns the archive directly; a near match on \struct{StructuralHash} (the content hash with extension-point body hashes stripped) warm-starts from a structurally similar archive.
Body-similar archives reduce convergence by approximately five-to-ten times for research variants of a familiar template.

\textbf{Federation.}
The archive is federation-shareable at Cipher's cold tier.
Two installations running the same model produce comparable archives; sharing them via a shared cold-tier object store enables a kernel-compilation genome that grows monotonically across the user community.
The privacy guarantees follow from \IR{2}'s vendor-neutrality: an L1 KernelCache entry contains kernel templates, recipes, and tile specs; no weights, no training data, no proprietary identifiers.

\subsection{Calibration: how the simulator reaches 95-to-98-percent accuracy}
\label{sec:mimic-calibration}

The simulator's accuracy target is 95-to-98-percent on Forge-emitted instruction streams.
Without calibration, the simulator would be 70-to-80-percent accurate (ballpark predictions from hardware specifications).
With calibration, the additional fifteen-to-twenty percentage points come from a structured microbenchmark suite that fits the per-chip constants.

The calibration harness has seven stages.

\textbf{Stage 1: Isolation.}
Single-warp, single-pipe, no memory access.
Approximately thirty microbenchmarks per chip target measure per-opcode latency, occupancy, and throughput.

\textbf{Stage 2: Functional-unit contention.}
Multiple warps doing the same operation.
Approximately twenty microbenchmarks measure per-pipe instance count and partition share.

\textbf{Stage 3: Memory isolation.}
Single-warp pointer-chase with varying stride.
Approximately fifteen microbenchmarks measure L1, L2, and main-memory latencies in isolation.

\textbf{Stage 4: Memory bandwidth.}
Maximum-width strided loads with all warps.
Approximately ten microbenchmarks measure aggregate-bandwidth ceilings at each level.

\textbf{Stage 5: Scheduler.}
Two-to-four warps with asymmetric stall patterns.
Approximately ten microbenchmarks fit the GTO greedy probability, the boost threshold, and the burst limit.

\textbf{Stage 6: Integration.}
Tiled GEMM, softmax, FlashAttention.
Approximately twenty microbenchmarks validate that the per-mechanism calibrations compose correctly.

\textbf{Stage 7: Regression.}
Five hundred real kernels with measured performance.
The accurate-tier residual is required to satisfy a P95 cap below five percent.

The total per-chip calibration time is approximately one to two months of engineering effort to reach 95-percent accuracy; reaching 98 percent requires another two-to-three months and provides diminishing returns for machine-learning workloads.

\subsection{Per-vendor backends}
\label{sec:mimic-per-vendor}

The per-vendor backends share the template (Section~\ref{sec:mimic-template}) and the public interface (Section~\ref{sec:mimic-public-interface}); their internals diverge by vendor.
The architecture-level overview of each follows; the misc/MIMIC.md design specification documents the vendor-specific instruction encoders, calibrated tables, simulator specializations, runtime ioctl protocols, and collective implementations.

\textbf{NVIDIA (\code{mimic/nv/}).}
\IR{3}NV is the Mercury SASS instruction stream for sm\_90 (Hopper), sm\_100 (Blackwell datacenter), sm\_103 (Blackwell Ultra), sm\_110 (Jetson Thor), sm\_120 and sm\_121 (Blackwell consumer and DGX Spark).
The encoder produces 128-bit Mercury instruction words and packages them into cubin ELF.
The simulator includes the Hopper warp scheduler (GTO with boost), the WGMMA and TMA async subsystems, the Blackwell tcgen05 with TMEM, and the four-level memory hierarchy.
The runtime library wraps \code{/dev/nvidia*} ioctls, replacing libcuda, libcudart, libnvJitLink, and libnvrtc.
The collective library implements ring, tree, halving-doubling, and hierarchical algorithms over NVLink (intra-node) and RDMA (inter-node).
NVIDIA is Mimic's first and reference-complete backend.

\textbf{AMD (\code{mimic/am/}).}
\IR{3}AM is AMDGPU ISA for CDNA3 (gfx942 MI300X, gfx950 MI325X) and RDNA3 and later for consumer.
The encoder may reuse LLVM-AMDGPU internally for ISA selection (the generator tables are open); the HSACO wrapper and the scheduler are Mimic's own.
The runtime library wraps \code{/dev/kfd} and per-device \code{/dev/dri/renderD*}.
The collective library targets XGMI intra-node and RoCE inter-node.

\textbf{Google TPU (\code{mimic/tpu/}).}
\IR{3}TPU targets v5p, v5e, v6e, and v7.
The encoder produces TPU executable format (protobuf with opaque payload).
The runtime library wraps \code{/dev/accel*}.
The collective library implements the ICI torus topology with native in-fabric aggregation.

\textbf{AWS Trainium (\code{mimic/trn/}).}
\IR{3}TRN targets trn1, trn2, and trn3.
The encoder produces NEFF (Neuron Executable File Format).
The runtime library wraps \code{/dev/neuron*}.
The collective library composes NeuronLink intra-node with EFA inter-node.

\textbf{Cerebras (\code{mimic/cer/}).}
\IR{3}CER targets WSE3 and later.
The encoder produces Cerebras Software Language binary format.
The runtime library wraps \code{/dev/cerebras*}.
The collective library targets the SwarmX wafer fabric.

\textbf{CPU (\code{mimic/cpu/}).}
\IR{3}CPU targets x86\_64 (AVX-512 and AVX2), AArch64 (NEON and SVE2), and RISC-V (V extension).
The encoder writes a native ELF shared object or stages C\nolinebreak\hspace{0.05em}\raisebox{0.4ex}{\tiny ++}\, source plus invokes a JIT compiler.
The runtime library uses standard \code{malloc}, \code{mmap}, and pthreads.
The CPU backend is the correctness oracle for the cross-vendor numerics CI: every \code{BITEXACT\_STRICT} recipe is realized as a scalar FMA chain on CPU, and every other backend's output is compared against the CPU reference.

\subsection{The runtime library replacement}
\label{sec:mimic-runtime}

Each per-vendor backend ships its own runtime library at \code{mimic/<vendor>/rt/} replacing libcuda, libnrt, libtpu, libhsa, with the same surface and zero vendor-SDK dependency.
The runtime library provides device discovery and open (one ioctl), memory-pool management (one large allocation at Keeper init, then position-independent pointer arithmetic), binary loading (the per-vendor binary format is parsed directly without vendor-supplied loaders), kernel submission (via the per-vendor pushbuffer or doorbell or DMA mechanism), graph capture (where supported by the vendor), and synchronization and completion polling.

The submission path is the latency-critical surface.
Each per-vendor backend documents the byte sequence the runtime emits per kernel launch.
For NVIDIA Hopper: a four-dword pushbuffer entry plus a doorbell write, totaling thirty-two bytes of host-to-GPU traffic per launch with a CPU critical path of approximately 80 to 200 nanoseconds.
For AMD CDNA3: a PM4 packet sequence plus a doorbell.
For TPU and Trainium: a DMA-descriptor write to the device.

The runtime library architecture and the per-vendor latency budgets are documented in detail in CRUCIBLE.md and MIMIC.md.

\subsection{The collective library replacement}
\label{sec:mimic-collective}

Each per-vendor backend ships its own collective library at \code{mimic/<vendor>/comm/} replacing NCCL, RCCL, libnccom, hcoll, with the same surface and zero vendor-collective-library dependency.
The collective library implements the standard primitives (all-reduce, all-gather, reduce-scatter, all-to-all, broadcast, point-to-point) with deterministic-reduction discipline (canonical pairwise-tree ordering by sorted UUID under \code{BITEXACT} recipes; algorithm selection per message size under \code{ORDERED}; opportunistic under \code{UNORDERED}).

A distinguishing structural property: in Mimic, collectives are \emph{kernels} in the \struct{ExecutionPlan} (Chapter~\ref{sec:execution-plan}), not separate runtime entities driven by a proxy thread.
A ring all-reduce is approximately two hundred bytes of pushbuffer (eight hops at sixteen bytes plus seven barriers), submitted in the same doorbell as the surrounding compute kernels.
NCCL by contrast issues per-collective \code{cuLaunchKernel}, costing approximately four-to-six microseconds of host overhead per collective; the proxy thread that drives NCCL's progress is a separate runtime construct.
The cost difference is structural, not engineering: collectives in Mimic move at the speed of pushbuffer and the fabric.

\subsection{Properties}
\label{sec:mimic-properties}

The properties below are commitments of the Mimic specification.
Each is enforced either structurally (by the per-vendor backend's type system) or by the cross-vendor numerics CI matrix (Chapter~\ref{sec:cross-vendor-ci}).

\begin{property}[Vendor-SDK independence]
\label{prop:mimic-no-sdk}
No per-vendor backend links against any vendor SDK runtime library.
Each backend wraps the kernel-driver ioctl protocol directly.
The kernel driver is the only vendor-proprietary surface Crucible inherits.
\end{property}

\begin{property}[Per-vendor isolation]
\label{prop:mimic-isolation}
The per-vendor backends are mutually independent.
Adding a new vendor changes no other backend, no shared-core code, and no Forge code.
The five public entry points are uniform; the dispatch is mechanical.
\end{property}

\begin{property}[Recipe-realization fidelity]
\label{prop:mimic-recipe-fidelity}
Each per-vendor backend implements \code{realize\_recipe}, mapping a \struct{NumericalRecipe} to vendor-specific instructions.
The realization is deterministic, documented per backend, and verified by the cross-vendor numerics CI matrix.
A backend that violates the declared cross-vendor tolerance fails the build.
\end{property}

\begin{property}[Calibrated cost-model fidelity]
\label{prop:mimic-calibrated}
Each per-vendor simulator's medium-tier output satisfies a P95 absolute-error cap below five percent against measured hardware on the per-vendor regression set of approximately five hundred kernels.
The accuracy target is preserved by Augur's continuous monitoring: drift triggers \code{mimic::recalibrate}, which re-runs the calibration harness during a cluster idle window.
\end{property}

\begin{property}[Round-trip-correct encoder]
\label{prop:mimic-roundtrip}
Each per-vendor backend's encoder satisfies \code{decode(encode(prog))} equals \code{prog} for every valid \IRstar{3} program.
The property is asserted by continuous integration tests that generate random programs, encode and decode, and compare structurally.
\end{property}
